\documentclass[a4paper,11pt]{article}
\usepackage{amsmath,amsthm,amsfonts,amssymb,amscd,amstext,vmargin,graphics,graphicx,tabularx,multicol} 
\usepackage[francais]{babel}
\usepackage[utf8]{inputenc}  
\usepackage[T1]{fontenc} 
\usepackage{pstricks-add,tikz,tkz-tab,variations}
\usepackage[autolanguage,np]{numprint} 
\usepackage{calc}

\setmarginsrb{1.5cm}{0.5cm}{1cm}{0.5cm}{0cm}{0cm}{0cm}{0cm} %Gauche, haut, droite, haut
\newcounter{numexo}
\newcommand{\exo}[1]{\stepcounter{numexo}\noindent{\bf Exercice~\thenumexo} : \marginpar{\hfill /#1}}
\reversemarginpar

\newcommand{\bmul}[1]{\begin{multicols}{#1}}
\newcommand{\emul}{\end{multicols}}

\newcounter{enumtabi}
\newcounter{enumtaba}
\newcommand{\q}{\stepcounter{enumtabi} \theenumtabi.  }
\newcommand{\qa}{\stepcounter{enumtaba} (\alph{enumtaba}) }
\newcommand{\initq}{\setcounter{enumtabi}{0}}
\newcommand{\initqa}{\setcounter{enumtaba}{0}}

\newcommand{\be}{\begin{enumerate}}
\newcommand{\ee}{\end{enumerate}}
\newcommand{\bi}{\begin{itemize}}
\newcommand{\ei}{\end{itemize}}
\newcommand{\bp}{\begin{pspicture*}}
\newcommand{\ep}{\end{pspicture*}}
\newcommand{\bt}{\begin{tabular}}
\newcommand{\et}{\end{tabular}}
\renewcommand{\tabularxcolumn}[1]{>{\centering}m{#1}} %(colonne m{} centrée, au lieu de p par défault) 
\newcommand{\tnl}{\tabularnewline}

\newcommand{\trait}{\noindent \rule{\linewidth}{0.2mm}}
\newcommand{\hs}[1]{\hspace{#1}}
\newcommand{\vs}[1]{\vspace{#1}}

\newcommand{\N}{\mathbb{N}}
\newcommand{\Z}{\mathbb{Z}}
\newcommand{\R}{\mathbb{R}}
\newcommand{\C}{\mathbb{C}}
\newcommand{\Dcal}{\mathcal{D}}
\newcommand{\Ccal}{\mathcal{C}}
\newcommand{\mc}{\mathcal}

\newcommand{\vect}[1]{\overrightarrow{#1}}
\newcommand{\ds}{\displaystyle}
\newcommand{\eq}{\quad \Leftrightarrow \quad}
\newcommand{\vecti}{\vec{\imath}}
\newcommand{\vectj}{\vec{\jmath}}
\newcommand{\Oij}{(O;\vec{\imath}, \vec{\jmath})}
\newcommand{\OIJ}{(O;I,J)}


\newcommand{\reponse}[1][1]{%
\multido{}{#1}{\makebox[\linewidth]{\rule[0pt]{0pt}{20pt}\dotfill}
}}

\newcommand{\titre}[5] 
% #1: titre #2: haut gauche #3: bas gauche #4: haut droite #5: bas droite
{
\noindent #2 \hfill #4 \\
#3 \hfill #5

\vspace{-1.6cm}

\begin{center}\rule{6cm}{0.5mm}\end{center}
\vspace{0.2cm}
\begin{center}{\large{\textbf{#1}}}\end{center}
\begin{center}\rule{6cm}{0.5mm}\end{center}
}



\begin{document}
\pagestyle{empty}
\titre{TP avec Geogebra : Symétrie axiale}{}{}{6ème}{}
\vspace*{0.15cm}

Ouvrir le logiciel Geogebra.\\
Préparer la feuille en enlevant les axes. \textit{(clic droit de la souris puis axes)}\\

\textbf{PARTIE 1 :} \textit{Symétrique d'un point par rapport à une droite}\\
\q Construire une droite (AB) et placer un point M qui n'appartient pas à la droite (AB). \\
Colorier la droite (AB) en rouge. \textit{(outil déplacer \includegraphics[scale=1]{curseur.eps}  , clic droit de la souris, propriétés, couleur)}\\
\q Construire le symétrique M' du point M par rapport à la droite (AB).\textit{ (outil symétrie axiale \includegraphics[scale=0.8]{outilsymax.eps}  )}\\
\q Tracer le segment [MM'], il coupe la droite (AB) en I, placer le point I.\\
 \textit{(outil intersection entre deux objets \includegraphics[scale=1]{intersection.eps}  )}\\
\q Mesurer les segments [MI] et [IM'] avec un outil de mesure de Geogebra. \\
\textbf{Le point I est . . . . . . . . . . . . . . . . . . . . . du segment [MM'].}\\
\q Faire afficher la mesure de l'angle $\widehat{MIA}$. \textit{(outil angle \includegraphics[scale=1]{outilangle.eps} )}\\
Que peut-on dire de la droite (AB) pour le segment [MM’] ?\\
   \textbf{La droite (AB) est . . . . . . . . . . . . . . . . . . . au segment [MM’].}\\

$\rightarrow$ \textbf{Enregistrez votre fichier} dans le dossier devoir sous le nom "TP-prenom "\\

\vspace*{0.15cm}

\textbf{PARTIE 2 :} \textit{Symétrique d'un segment par rapport à une droite}\\
 \initq \q Sur le même fichier, se déplacer si besoin avec l'outil \includegraphics[scale=0.85]{outilsdeplacement.eps}  et construire un segment [CD].\\
\q Construire le symétrique [C’D’] du segment [CD] par rapport à la droite (AB).\\
\q Faire afficher et comparer les longueurs des segments [CD] et [C’D’].\\

\noindent\fbox{\parbox{\linewidth-2\fboxrule-2\fboxsep}{\textbf{\underline{Conclusion }:} \textit{Le symétrique d'un segment par rapport à une droite est} . . . . . . . . . . . . . . . . . . . . . . . . . . . . . . . . . . . . . . .}}\\

\vspace*{0.15cm}

\textbf{PARTIE 3 :} \textit{Symétrique d'une droite par rapport à une droite}\\
 \initq \q Toujours sur le même fichier, se déplacer si besoin avec l'outil \includegraphics[scale=0.85]{outilsdeplacement.eps}  et construire une droite (EF).\\
\q Puis construire son symétrique par rapport à la droite (AB).\\


\noindent\fbox{\parbox{\linewidth-2\fboxrule-2\fboxsep}{\textbf{\underline{Conclusion }:} \textit{Le symétrique d'une droite par rapport à une droite est} . . . . . . . . . . . . . . . . . . . . . . . . . . . . . }}\\

$\rightarrow$ \textbf{Penser à enregistrer} régulièrement votre travail en cliquant sur "Sauvegarder".\\

\vspace*{0.15cm}

\textbf{PARTIE 4 :} \textit{Symétrique d'un cercle par rapport à une droite}\\
 \initq \q Toujours sur le même fichier, se déplacer si besoin avec l'outil \includegraphics[scale=0.85]{outilsdeplacement.eps}  et construire un cercle de centre O et passant par un point G.\\
\q Construire ensuite les symétriques O’ et G’ de O et G par rapport à la droite (AB) puis le symétrique du cercle.\\
\q Tracer les segments [OG] et [O’G’] et afficher leurs longueurs.\\


\noindent\fbox{\parbox{\linewidth-2\fboxrule-2\fboxsep}{\textbf{\underline{Conclusion }:} \textit{Le symétrique d'un cercle  par rapport à une droite est} . . . . . . . . . . . . . de même . . . . . . . . . }}\\
\vspace*{0.15cm}

\textbf{PARTIE 5 :} \textit{Symétrique d'un polygone par rapport à une droite}\\
\initq \q Toujours sur le même fichier, construire un polygone RSTUV puis son symétrique R’S’T’U’V’ par rapport à la droite (AB).\\
\q Afficher et comparer les longueurs des segments [RS], [ST], [TU], [UV], [VR] et leurs symétriques.\\
\q Afficher et comparer les mesures de l'angle $\widehat{RST}$ et de son symétrique.\\
\q Afficher et comparer l'aire du polygone RSTUV et de son symétrique.\\


\noindent\fbox{\parbox{\linewidth-2\fboxrule-2\fboxsep}{\textbf{\underline{Conclusion }:} La symétrie axiale conserve : - . . . . . . . . . . . . . . . . . .\\
\hspace*{7.5cm}- . . . . . . . . . . . . . . . . . .\\
\hspace*{7.5cm}- . . . . . . . . . . . . . . . . . .\\
} 
}\\
\end{document}
